\documentclass[compress,t]{beamer}
% \documentclass[compress,notes]{beamer}
% \documentclass[compress,notes=only]{beamer}

\newif\iffinal
\finalfalse

\usetheme{Paul}
% \setbeamercovered{transparent}
\usepackage[utf8]{inputenc}
\usepackage[ngerman]{babel}
\usepackage{floatflt}
\usepackage{graphicx}
\usepackage{tikz}
\usepackage{xcolor}
\usepackage{listings}
\usetikzlibrary{shapes,arrows,calc,decorations.markings}

\usepackage{mdframed}

\newcommand{\qq}[1]{\symbol{34}#1\symbol{34}}
\newcommand{\bash}[1]{
    \begin{mdframed}[backgroundcolor=black!80,
            innerbottommargin=1pt,
            innertopmargin=3pt,
            innerleftmargin=3pt,
            innerrightmargin=3pt,
            linewidth=0]
        \textcolor{white}{\texttt{\small #1}}
    \end{mdframed}
}

\title{Git}
\author{Paul Bienkowski (2bienkow)}
\institute{KunterBuntesSeminar}
\date{2013-05-29}
\setbeamerfont{note page}{size=\small}

\tikzset{
    onslide/.code args={<#1>#2}{% http://tex.stackexchange.com/a/6155/16595
        \only<#1>{\pgfkeysalso{#2}}
    }
}

\tikzset{
    commit/.style={
        draw,
        fill=cyan!30,
        ellipse,
        font=\tiny,
        minimum height=1.5em,
        text centered,
        minimum width=2em,
    },
    branch/.style={
        draw,
        fill=brown!30,
        rectangle,
        rounded corners=1pt,
        font=\tiny,
        minimum height=1.1em,
        text centered,
        minimum width=2em,
    },
    current/.style={
        fill=orange!50
    },
    child/.style={
        latex-,
        thick
    },
    pointer/.style={
        pointer'
    },
    pointer'/.style={
        draw=gray,
        thick,
        densely dotted
    },
    repo/.style={
        draw,
        fill=cyan!30,
        circle,
        font=\tiny,
        minimum height=3em,
        minimum width=3em,
        text centered,
    },
    remote/.style={
        decoration={
            markings,
            mark={
                at position 0
                with {
                    \fill circle [xshift=1pt,radius=1pt];
                }
            }
        },
        postaction=decorate,
        -latex,
        thick
    }
}

%\AtBeginSection[]{
    %\frame{
    %\frametitle{Outline}
    %\tableofcontents[current]
    %}
%}

\begin{document}
\frame{
    \titlepage
}

\frame{
    \frametitle{Gliederung}
    \tableofcontents
}

\section{Einleitung}
\subsection{Wozu Git?}
\frame {
    \frametitle{Wozu Git?}

    \uncover<0->{Lasst uns zusammen ein Programm schreiben!}
    \uncover<2->{\textit{(Evolution eines Hackers)}}

    \begin{enumerate}
        \item<3-> Wir arbeiten im Etherpad!
        \item<4-> Wir arbeiten auf \texttt{rzssh1}!
        \item<5-> Du schickst mir dann die neue Version per Mail!
        \item<6-> Du schickst mir dann ein Diff per Mail!
        \item<7-> Wir nehmen Versionskontrolle!
    \end{enumerate}

    \uncover<8> {
        $\rightarrow$ gilt auch für Hausaufgaben in \LaTeX
    }

    \vspace{0.2cm}

    \makebox[\textwidth][c]{
        \begin{tikzpicture}[x=1cm,y=1cm]
            \uncover<3->{
                \node[anchor=south west,inner sep=0,opacity=0.2] (human-01) at (0,0) {\includegraphics[scale=0.5]{gfx/human-01.png}};
            }
            \uncover<4->{
                \node[anchor=south west,inner sep=0,opacity=0.2] (human-02) at (1.2,0) {\includegraphics[scale=0.5]{gfx/human-02.png}};
            }
            \uncover<5->{
                \node[anchor=south west,inner sep=0,opacity=0.2] (human-03) at (2.1,0) {\includegraphics[scale=0.5]{gfx/human-03.png}};
            }
            \uncover<6->{
                \node[anchor=south west,inner sep=0,opacity=0.2] (human-04) at (2.9,0) {\includegraphics[scale=0.5]{gfx/human-04.png}};
            }
            \uncover<7->{
                \node[anchor=south west,inner sep=0,opacity=0.2] (human-05) at (4.0,0) {\includegraphics[scale=0.5]{gfx/human-05.png}};
            }
        \end{tikzpicture}
    }
}

\frame{
    \frametitle{Was ist Versionierung?}

    \begin{itemize}
        \item \emph{Snapshots} einzelner Dateiversionen speichern
        \item History aufbewahren und wiederherstellbar machen
        \item verschiedene Versionen zusammenführen
        \item \emph{(optional)} Synchronisation mit entfernten Kopien (Kollaboration)
    \end{itemize}
    $\rightarrow$ mehr als nur ein Backup
}

\subsection{Warum gerade git?}
\frame{
    \frametitle{Warum gerade git?}

    Vorteile:
    \begin{itemize}
        \item verteilt (serverunabhängig, jeder \emph{clone} ist eigenständig/vollständig)
        \item schnell (lokal, Implementation in C)
        \item optimal für Quelltext
        \item FOSS
    \end{itemize}

    Nachteile:
    \begin{itemize}
        \item wenig geeignet für Binärdateien
        \item gesamte History kann groß werden
    \end{itemize}
}

\section{Git Internals}
\subsection{Objekte}
\frame{
    \frametitle{Objekte}
    \begin{itemize}
        \item Git als map-type storage
        \item Hashes bilden auf Daten ab
        \item Speicherung in einzelnen Dateien (``der Kernel macht das'')
        \item jedes Objekt in git ist eine Datei
            \begin{itemize}
                \item Versions-Snapshot von Dateien (\emph{blob})
                \item Versions-Snapshot von Verzeichnissen (\emph{tree})
                \item Commit-Informationen (\emph{commit})
            \end{itemize}
        \item \emph{tree} und \emph{commit} enthalten Informationen in Klartext,  \emph{blob} den Dateiinhalt
    \end{itemize}

    \pause
    \vspace{0.5em}
    \makebox[\textwidth][c]{
        \begin{tikzpicture}[x=1cm,y=1cm]
            \node (rect) at (0, 0) [draw,thick,minimum width=2cm,minimum height=0.6cm] {type};
            \node (rect) at (2, 0) [draw,thick,minimum width=2cm,minimum height=0.6cm] {length};
            \node (rect) at (5, 0) [draw,thick,minimum width=4cm,minimum height=0.6cm] {data};
            \node (rect) at (3, -0.1) [draw,thick,minimum width=8.2cm,minimum height=1cm,label=below:\emph{zlib compress}] {};
        \end{tikzpicture}
    }
}

\subsection{Remotes}
\frame {
    \frametitle{Clones}

    \begin{itemize}
        \item in einem Repository entstehen nur neue Objekte (Dateien), sie werden
            nicht verändert
        \item jede Datei hat einen eindeutigen Namen (SHA kollidiert
            \textit{praktisch} nicht)
        \item $\rightarrow$ Synchronisation ohne zentrale Verwaltung möglich
    \end{itemize}

    Einzige Ausnahme: Branch-Referenzen verändern sich.
}

\frame {
    \frametitle{Remotes}

    \vspace{0.5em}
    \makebox[\textwidth][c]{
        \begin{tikzpicture}[x=1cm,y=1cm]
            \node[repo] (github) at (0, 2) {\textit{github}};
            \node[repo] (paul) at (-1, 0) {\textit{paul}};
            \node[repo] (rzssh1) at (1, 0) {\textit{rzssh1}};

            \draw[remote] (paul) -- node[midway, left]
                {\tiny \textbf{origin} $\rightarrow$ git@github.com:opatut/repo.git} (github);
            \draw[remote] (rzssh1) -- node[midway, right]
                {\tiny \textbf{origin} $\rightarrow$ git@github.com:opatut/repo.git} (github);

            %\draw[remote] (paul) -- node[midway, below] {\tiny{}ssh://2bienkow@rzssh1.informatik.uni-hamburg.de:/home/scratch/2bienkow/repo} (rzssh1);
        \end{tikzpicture}
    }
}

\section{Basics}

\subsection{Working copy}
\frame{
    \frametitle{Working copy}
    \begin{itemize}
        \item In \texttt{.git/} liegen alle Versionen aller Dateien als \emph{blob} vor
        \item Wie soll man damit arbeiten?
        \item $\rightarrow$ aktuelle Version (\emph{HEAD}) liegt im Hauptverzeichnis
        \item verständliche Dateinamen (statt Hashes)
        \item wechseln der Version per \textsc{checkout}
            \bash{git checkout 4b5c8e2f95a4407c4d0c596565b367eaca07af57}
            $\rightarrow$ nicht möglich, wenn unversionierte Änderungen vorliegen
    \end{itemize}

}

\subsection{Branching}
\frame{
    \frametitle{Branching}

    \begin{itemize}
        \item Commits liegen ungeordnet vor
        \item $\rightarrow$ Welche ist die ``aktuelle'' Version?
        \item Ein \textbf{Branch} kann auf einen Commit zeigen (``Pointer'')
        \item Branches haben \textbf{Namen} (z.B. \textit{master} oder \textit{my-feature})
        \item Man kann zwischen Branches wechseln wie zwischen Commits (\textsc{checkout})
        \item Nach dem committen zeigt der aktuelle Branch auf den neuen Commit
    \end{itemize}
}

\frame{
    \frametitle{Branching}

    \only<1>{\bash{git status}}
    \only<2>{\bash{git checkout -b foobar}}
    \only<3>{\bash{git commit -m \qq{D}}}
    \only<4>{\bash{git checkout master}}
    \only<5>{\bash{git checkout -b hotfix}}
    \only<6>{\bash{git commit -m \qq{E}}}
    \only<7>{\bash{git checkout master}}
    \only<8>{\bash{git merge hotfix}}
    \only<9>{\bash{git branch -d hotfix}}
    \only<10>{\bash{git checkout foobar}}
    \only<11-12>{\bash{git commit -m \qq{F}}}
    \only<13>{\bash{git checkout master}}
    \only<14>{\bash{git merge foobar}}
    \only<15>{\bash{git branch -d foobar}}

    \vspace{1em}

\iffinal
    \makebox[\textwidth][c]{
        \begin{tikzpicture}[x=3em,y=2em]
            \uncover<1-> {
                \node[commit] (a) at (0, 0) {A};
                \node[commit] (b) at (1, 0) {B};
                \node[commit] (c) at (2, 0) {C};

                \draw[child] (a) -> (b);
                \draw[child] (b) -> (c);
            }

            \uncover<1-7> {
                \node[branch,onslide=<1>current,onslide=<4>current,onslide=<7>current] (master) at (2, 1) {master};
                \draw[pointer] (master) -> (c);
            }

            \uncover<2> {
                \node[branch,current] (foobar) at (2, -1) {foobar};
                \draw[pointer] (foobar) -> (c);
            }

            \uncover<3-> {
                \node[commit] (d) at (3, -1) {D};
                \draw[child,onslide=<12-13>{red}] (c) -> (d);
            }

            \uncover<3-10> {
                \node[branch,onslide=<3>current,onslide=<10>current] (foobar) at (3, -2) {foobar};
                \draw[pointer] (foobar) -> (d);
            }

            \uncover<5> {
                \node[branch,current] (hotfix) at (2, 2) {hotfix};
                \draw[pointer'] (hotfix) -- (master);
            }

            \uncover<6-> {
                \node[commit] (e) at (3, 0) {E};
                \draw[child,onslide=<12-13>{green}] (c) -> (e);
            }

            \uncover<6-8> {
                \node[branch,onslide=<6>{current}] (hotfix) at (3, 1) {hotfix};
                \draw[pointer] (hotfix) -> (e);
            }

            \uncover<8> {
                \node[branch,current] (master) at (3, 2) {master};
                \draw[pointer'] (master) -> (hotfix);
            }

            \uncover<9-13> {
                \node[branch,onslide=<9>current,onslide=<13>current] (master) at (3, 1) {master};
                \draw[pointer] (master) -> (e);
            }


            \uncover<11-> {
                \node[commit] (f) at (4, -1) {F};
                \draw[child,onslide=<12-13>{red}] (d) -> (f);
            }

            \uncover<11-14> {
                \node[branch,onslide=<-12>current] (foobar) at (4, -2) {foobar};
                \draw[pointer] (foobar) -> (f);
            }

            \uncover<14-> {
                \node[commit] (g) at (5, 0) {G*};
                \draw[child] (e) -> (g);
                \draw[child] (f) -> (g);
            }

            \uncover<14-> {
                \node[branch,current] (master) at (5, 1) {master};
                \draw[pointer] (master) -> (g);
            }

            \uncover<15> {}
        \end{tikzpicture}
    }
\else
    -- tikz picture here, finaltrue to show --
\fi
}

\subsection{Merging}
\frame{
    \frametitle{Merging}

    Was passiert?
    \begin{itemize}
        \item gemeinsamen Vorgänger finden
        \item Änderungen ermitteln
        \item beide Änderungssätze auf gemeinsamen Vorgänger anwenden
        \item neuen Commit erstellen (automatische \emph{message})
    \end{itemize}

    Konflikte
    \begin{itemize}
        \item bei Änderung gleicher Zeile kann git nicht entscheiden, welche
            Änderung übernommen werden soll
        \item
            {\scriptsize
            \texttt{Auto-merging <filename>}\\
            \texttt{CONFLICT (content): Merge conflict in <filename>}\\
            \texttt{Automatic merge failed; fix conflicts and then commit the result.}\\
            }
    \end{itemize}
}

\begin{frame}[fragile]
    \frametitle{Merging - Konflikte}

    \begin{verbatim}
vor dem Konflikt
<<<<<<< HEAD
meine Änderungen
=======
deine Änderungen
>>>>>>> fremder-branch
nach dem Konflikt
    \end{verbatim}

\end{frame}

\end{document}
